% Generated draft: requires mathematical and source review.
\documentclass[11pt]{article}
\usepackage[margin=25mm]{geometry}
\usepackage{fontspec}
\setmainfont{Latin Modern Roman}
\usepackage{amsmath,amssymb,mathtools,unicode-math}
\setmathfont{Latin Modern Math}
\usepackage{enumitem,needspace,xurl,hyperref}
\setlength{\parindent}{0pt}
\setlength{\parskip}{5pt}
\setlength{\emergencystretch}{3em}
\newcommand{\N}{\mathbb{N}}
\newcommand{\Z}{\mathbb{Z}}
\newcommand{\Q}{\mathbb{Q}}
\newcommand{\R}{\mathbb{R}}
\newcommand{\C}{\mathbb{C}}
\newcommand{\F}{\mathbb{F}}
\newcommand{\A}{\mathbb{A}}
\newcommand{\PP}{\mathbb{P}}
\newcommand{\E}{\mathbb{E}}
\newcommand{\eps}{\varepsilon}
\newcommand{\dd}{\,\mathrm{d}}
\begin{document}
\section*{UnsolvedMath ID: 30006369}
\textbf{Source problem code:} catalogue entry
\subsection{Definitions and mathematical statement}
For terminology, let \((M,\omega)\) be a symplectic manifold and let \(\phi\colon M\to M\) be a Hamiltonian diffeomorphism. A point \(x\in M\) is periodic if \(\phi^k(x)=x\) for some integer \(k\geq 1\); its least period is the smallest such integer. Thus
\[
\operatorname{Per}(\phi)=\bigcup_{k\geq 1}\{x\in M:\phi^k(x)=x\}.
\]
The supplied background places the question in an asymptotically linear Hamiltonian setting, but gives neither the ambient symplectic manifold nor the precise asymptotic hypotheses. Consequently, the intended class of maps cannot be specified from this record, and the question must not be read as asserting a conjecture for every Hamiltonian diffeomorphism on every symplectic manifold.

Let \(\operatorname{Per}_{\mathrm{nt}}(\phi)\subseteq\operatorname{Per}(\phi)\) denote the points called ``non-trivial'' in the source, and let \(h(\phi)\) denote the source's polynomial growth exponent for these periodic points as their period increases. These are placeholders for source definitions that are absent from the supplied record: in particular, ``non-trivial'' is not defined, and no counting function or limiting formula for \(h(\phi)\) is supplied.

With these unresolved qualifications, the conjectural statement is that every map \(\phi\) in the intended asymptotically linear Hamiltonian class satisfies
\[
\#\operatorname{Per}_{\mathrm{nt}}(\phi)=\infty
\quad\Longrightarrow\quad
h(\phi)\geq 2.
\]
The source describes this conclusion as growth that is ``at least quadratic.'' Without the definition of \(h\), no particular counting-function inequality can be inferred from that description. This is a conditional statement template requiring recovery of the missing hypotheses and definitions before it becomes a precise, self-contained conjecture.
\subsection{Short English statement}
The supplied conjecture asks whether infinitely many non-trivial periodic points force a polynomial growth exponent of at least \(2\) in the intended asymptotically linear Hamiltonian setting. The admissible maps, the meaning of non-triviality, and the growth exponent require definitions absent from the record.
\subsection{Sources}
\begin{itemize}
\item Local UnsolvedMath registry: \texttt{data/unsolvedmath/problems.json}, permanent ID \(30006369\), title ``Quadratic Growth of Periodic Points for Hamiltonian Maps,'' exact source problem code \texttt{catalogue entry}.
\item Supplied report citation: \emph{Dynamische Systeme (2026)}, Oberwolfach Reports, DOI \texttt{10.4171/owr/2025/31}, pp.\ 1647--1706. \url{https://doi.org/10.4171/owr/2025/31}. The record gives proposed year 2025 and no named proposer.
\item Supplied literature reference labeled ``Masci primary preprint'': \url{https://arxiv.org/abs/2403.01855}. No title, theorem locator, or explicit attribution of this conjecture to an individual is supplied.
\item The original excerpt attributes its inspiration to Franks' Poincar\'e--Birkhoff theorem, citing reference [6, Corollary 2.4]. The bibliographic identity of [6] is not supplied.
\item Literature assessment dated \texttt{2026-08-22}: the supplied curation labels the problem open and reports that the quadratic-growth assertion is presented as a conjecture, with no proof found in its review of the Masci preprint and the Oberwolfach report. This is supplied dated curation, not an independently verified assessment.
\item The registry labels the statement \texttt{corrected\_verified} and reports that the official report was checked and the notation repaired. These are supplied provenance claims; the sources have not been independently visited or verified for this draft.
\end{itemize}
\end{document}
